\section{Guia de estudos}

\subsection{Como estudar esta matéria}

Em vez de memorizar listas isoladas, acompanhe o caminho do material sedimentar:

\[
\boxed{
\begin{gathered}
\text{rocha-fonte}
\rightarrow
\text{intemperismo}
\rightarrow
\text{erosão e transporte}
\\
\rightarrow
\text{deposição}
\rightarrow
\text{diagênese}
\rightarrow
\text{rocha sedimentar}
\end{gathered}
}
\]

Para cada etapa, procure responder: \emph{qual é o processo, qual é o agente, qual
produto é gerado e que evidência fica preservada?}

\subsection{Resumo dos processos}

\begin{table}[H]
  \centering
  \caption{Relações fundamentais do sistema sedimentar.}
  \begin{tabularx}{\textwidth}{>{\bfseries}lXX}
    \toprule
    Etapa & Processo principal & Produto ou registro \\
    \midrule
    Intemperismo & Desagregação e decomposição da rocha & Clastos, argilominerais e íons dissolvidos \\
    Erosão & Remoção do material da área-fonte & Disponibilização de sedimento \\
    Transporte & Tração, saltação, suspensão e solução & Seleção, abrasão e orientação \\
    Deposição & Perda de competência/capacidade ou precipitação & Camadas e formas de leito \\
    Diagênese & Compactação, cimentação e reações pós-deposicionais & Litificação e modificação dos poros \\
    \bottomrule
  \end{tabularx}
\end{table}

\subsection{Distinções que não devem ser confundidas}

\begin{table}[H]
  \centering
  \caption{Conceitos próximos, mas diferentes.}
  \begin{tabularx}{\textwidth}{>{\bfseries}lXX}
    \toprule
    Par & Primeiro conceito & Segundo conceito \\
    \midrule
    Intemperismo/erosão & Altera a rocha no local & Remove o material \\
    Competência/capacidade & Maior tamanho transportável & Quantidade total transportável \\
    Arredondamento/esfericidade & Curvatura de arestas e vértices & Aproximação da forma a uma esfera \\
    Matriz/cimento & Material detrítico fino & Precipitado diagenético \\
    Porosidade/permeabilidade & Volume de vazios & Conectividade para o fluxo \\
    Bacia de sedimentação/bacia sedimentar & Local atual de deposição & Unidade geológica subsidente e preenchida \\
    Fácies/ambiente & Conjunto observável de características & Interpretação genética \\
    \bottomrule
  \end{tabularx}
\end{table}

\subsection{Equações e limites essenciais}

\begin{align*}
  \phi &= -\log_2 d_{\mathrm{mm}}, \\
  \mathrm{Fr} &= \frac{U}{\sqrt{gD}}.
\end{align*}

Interpretação do número de Froude:

\[
\mathrm{Fr}<1\;\text{subcrítico},
\qquad
\mathrm{Fr}=1\;\text{crítico},
\qquad
\mathrm{Fr}>1\;\text{supercrítico}.
\]

Limites granulométricos:

\[
\begin{array}{rcl}
d>\SI{2}{\milli\metre} &:& \text{cascalho},\\
\SI{0.0625}{\milli\metre}\le d\le\SI{2}{\milli\metre} &:& \text{areia},\\
\SI{0.0039}{\milli\metre}\le d<\SI{0.0625}{\milli\metre} &:& \text{silte},\\
d<\SI{0.0039}{\milli\metre} &:& \text{argila}.
\end{array}
\]

\subsection{Roteiro para descrever uma amostra siliciclástica}

Siga sempre a mesma sequência:

\begin{enumerate}
  \item determine a classe granulométrica dominante;
  \item estime a distribuição das frações cascalho, areia e lama;
  \item descreva o grau de seleção;
  \item descreva arredondamento e esfericidade separadamente;
  \item reconheça arcabouço, matriz, cimento e poros;
  \item identifique quartzo, feldspatos e fragmentos líticos;
  \item observe cor, estruturas e fósseis;
  \item atribua o nome da rocha;
  \item somente depois proponha processos e ambiente deposicional.
\end{enumerate}

\subsection{Classificação rápida das rochas siliciclásticas}

\begin{table}[H]
  \centering
  \caption{Critérios principais de classificação.}
  \begin{tabularx}{\textwidth}{>{\bfseries}lXX}
    \toprule
    Grupo & Critério principal & Subdivisões úteis \\
    \midrule
    Conglomerados & Clastos maiores que \SI{2}{\milli\metre} & Clasto/matriz-suportado; mono/oligo/polimítico \\
    Arenitos & Areia e composição Q--F--L & Arenita ou wacke; quartzosa, feldspática ou lítica \\
    Pelitos & Silte, argila e fissilidade & Siltito, argilito, lamito e folhelho \\
    \bottomrule
  \end{tabularx}
\end{table}

\subsection{Estruturas sedimentares e interpretação}

\begin{table}[H]
  \centering
  \caption{Estruturas e informações que podem fornecer.}
  \begin{tabularx}{\textwidth}{>{\bfseries}lX}
    \toprule
    Estrutura & Interpretação possível \\
    \midrule
    Estratificação cruzada & Migração de formas de leito e paleocorrente \\
    Ripple assimétrico & Corrente predominantemente unidirecional \\
    Ripple simétrico & Movimento oscilatório por ondas \\
    Gradação normal & Redução da energia durante a deposição \\
    Gretas de dessecação & Exposição subaérea e perda de água \\
    Marcas de sola & Erosão do substrato e possível sentido de fluxo \\
    Flaser/wavy/lenticular & Alternância de transporte de areia e decantação de lama \\
    Hummocky/swaley & Fluxo combinado relacionado a tempestades \\
    Imbricação & Orientação de clastos e possível paleocorrente \\
    Bioturbação & Atividade de organismos no substrato \\
    Estruturas de carga/escape & Instabilidade e movimentação de sedimento saturado \\
    \bottomrule
  \end{tabularx}
\end{table}

Uma estrutura isolada raramente determina um ambiente. A interpretação deve combinar
litologia, granulometria, geometria, contatos, paleocorrentes, fósseis e associação de
fácies.

\subsection{Perguntas de revisão}

\begin{enumerate}
  \item Qual é a diferença entre intemperismo físico e químico?
  \item Por que a água é importante no intemperismo químico?
  \item Que relação geral existe entre as séries de Bowen e Goldich?
  \item Diferencie erosão, transporte e deposição.
  \item Diferencie competência e capacidade de um fluxo.
  \item Quais são os quatro modos fundamentais de transporte sedimentar?
  \item Por que argilas exigem maior velocidade para serem erodidas do que alguns grãos de areia?
  \item O que significam \(\mathrm{Fr}<1\) e \(\mathrm{Fr}>1\)?
  \item Diferencie seleção, arredondamento e esfericidade.
  \item Quais são os limites fundamentais de cascalho, areia, silte e argila?
  \item Diferencie conglomerado clasto-suportado e matriz-suportado.
  \item Qual é a diferença entre arenita e wacke?
  \item O que representam os polos Q, F e L?
  \item Diferencie matriz e cimento.
  \item Quais processos fazem parte da diagênese?
  \item Como distinguir ripple de corrente e ripple de onda?
  \item O que uma estratificação espinha de peixe pode indicar?
  \item Por que fácies e ambiente deposicional não são sinônimos?
  \item Como soerguimento e subsidência participam do sistema sedimentar?
  \item O que uma discordância representa no registro geológico?
\end{enumerate}

\subsection{Respostas resumidas}

\begin{enumerate}
  \item O físico desagrega sem mudança mineralógica dominante; o químico transforma ou dissolve minerais.
  \item Ela atua como solvente, reagente e meio de transporte de íons.
  \item Minerais que cristalizam em temperaturas mais altas tendem a ser menos estáveis na superfície.
  \item Erosão remove, transporte desloca e deposição acumula.
  \item Competência refere-se ao tamanho máximo; capacidade, à quantidade total.
  \item Tração, saltação, suspensão e solução.
  \item Devido à coesão eletroquímica entre partículas finas.
  \item Fluxo subcrítico e supercrítico, respectivamente.
  \item Dispersão de tamanhos; curvatura das arestas; forma global.
  \item \SI{2}{\milli\metre}, \SI{0.0625}{\milli\metre} e \SI{0.0039}{\milli\metre}.
  \item No primeiro os clastos se tocam; no segundo ficam dispersos na matriz.
  \item A diferença principal é a proporção de matriz fina.
  \item Quartzo, feldspatos e fragmentos líticos.
  \item Matriz é detrítica; cimento é precipitado durante a diagênese.
  \item Compactação, cimentação, dissolução, recristalização e substituição.
  \item O de corrente tende a ser assimétrico; o de onda tende a ser simétrico.
  \item Reversão periódica das correntes, comum em sistemas de maré.
  \item Fácies é descrição; ambiente é interpretação.
  \item Soerguimento favorece exposição e erosão; subsidência cria acomodação.
  \item Um intervalo de não deposição e/ou erosão, isto é, tempo ausente.
\end{enumerate}
